TaperNorm replaces a pre-norm Transformer's internal normalization with a gated map that acts like RMSNorm early in training and then becomes a fixed, calibrated linear scaling. Without per-token normalization, such a model might forget more after a shift in the training data. We tested this in a small pilot with design, endpoint and decision rules fixed before training. Paired 17.7M-parameter models with internal RMSNorm or TaperNorm (three seeds) were trained on 150M web tokens, then continued for 100M tokens on web text or Python code. The primary endpoint is a difference-in-differences D in held-out web cross-entropy; positive D means extra forgetting under TaperNorm. Switching to Python raised web cross-entropy by about 1.66 nats/token in both models. Mean D was -0.022 nats/token (seeds -0.020, -0.021, -0.026), below the pre-specified +0.015 bound, so the pre-specified decision is to stop: we find no evidence that TaperNorm increases persistent forgetting in this setting. In exploratory analyses, D was near zero for most of continuation (after a brief early dip) and became negative as the learning rate annealed, and the code-specific activation-scale shift that motivated the hypothesis was small in both models. Code and all evaluation records are released.
